\documentclass[11pt,a4paper]{article}
\newcommand{\gradelabel}{Grade 9}
\newcommand{\leconhdr}{Simulation d'un feu intelligent}
% Préambule commun aux plans de cours de robotique (Grades 7 à 12).
% Chaque plan de cours commence par \documentclass[11pt,a4paper]{article}, définit
% \gradelabel et \leconhdr (en-têtes de page), puis % Préambule commun aux plans de cours de robotique (Grades 7 à 12).
% Chaque plan de cours commence par \documentclass[11pt,a4paper]{article}, définit
% \gradelabel et \leconhdr (en-têtes de page), puis % Préambule commun aux plans de cours de robotique (Grades 7 à 12).
% Chaque plan de cours commence par \documentclass[11pt,a4paper]{article}, définit
% \gradelabel et \leconhdr (en-têtes de page), puis \input{../../preambule_robotique.tex}.
\usepackage[T1]{fontenc}
\usepackage[utf8]{inputenc}
\usepackage{lmodern}
\usepackage[french]{babel}
\usepackage[margin=2cm,headheight=14pt]{geometry}
\usepackage[most]{tcolorbox}
\usepackage{xcolor,tabularx,booktabs,array,enumitem,fancyhdr,amssymb,tikz}
\usetikzlibrary{arrows.meta,positioning}

\pagestyle{fancy}
\fancyhf{}
\lhead{\small Robotique -- \gradelabel{} -- Novembre 2026}
\rhead{\small \leconhdr}
\cfoot{\thepage}
\setlist{itemsep=1pt,topsep=2pt,leftmargin=*}
\renewcommand{\arraystretch}{1.35}
\newcolumntype{Y}{>{\raggedright\arraybackslash}X}
\newcommand{\cb}{$\square$\,}
\newcommand{\fl}{$\rightarrow$}
\newcommand{\lig}{\noindent\rule{\linewidth}{0.3pt}\\[5mm]}

\newtcolorbox{etape}[3]{enhanced,breakable,colback=#1!6,colframe=#1!70!black,
  coltitle=white,fonttitle=\bfseries,title={#2\hfill #3},boxrule=0.8pt,arc=2mm,
  before skip=8pt,after skip=8pt}
\newtcolorbox{astuce}[1][Différenciation et besoins particuliers]{breakable,
  colback=yellow!12,colframe=yellow!50!black,fonttitle=\bfseries\small,
  title={#1},boxrule=0.5pt,arc=1mm,before skip=6pt,after skip=6pt}
\newtcolorbox{code}[1][Code]{breakable,colback=black!4,colframe=black!60,
  fonttitle=\bfseries\small,title={#1},boxrule=0.5pt,arc=1mm,fontupper=\ttfamily\small,
  before skip=6pt,after skip=6pt}
% Blocs de code coloré (écrire le code sans accents : comme dans l'éditeur)
\newtcblisting{arduino}[1][Code Arduino]{listing only,breakable,colback=black!4,
  colframe=black!60,fonttitle=\bfseries\small,title={#1},boxrule=0.5pt,arc=1mm,
  before skip=6pt,after skip=6pt,
  listing options={language=C++,basicstyle=\ttfamily\small,columns=fullflexible,
  keepspaces=true,showstringspaces=false,commentstyle=\color{green!40!black}\itshape,
  keywordstyle=\color{blue!70!black}\bfseries}}
\newtcblisting{python}[1][Code Python]{listing only,breakable,colback=black!4,
  colframe=black!60,fonttitle=\bfseries\small,title={#1},boxrule=0.5pt,arc=1mm,
  before skip=6pt,after skip=6pt,
  listing options={language=Python,basicstyle=\ttfamily\small,columns=fullflexible,
  keepspaces=true,showstringspaces=false,commentstyle=\color{green!40!black}\itshape,
  keywordstyle=\color{blue!70!black}\bfseries}}
\newcommand{\entete}[2]{%
  \begin{tcolorbox}[colback=#1!10,colframe=#1!70!black,boxrule=0.8pt,arc=2mm]
  {\Large\bfseries #2}\hfill Nom : \rule{4cm}{0.3pt}\quad Date : \rule{2.2cm}{0.3pt}
  \end{tcolorbox}}

% Bandeau de titre : #1 titre, #2 semaine
\newcommand{\titreblock}[2]{%
\begin{tcolorbox}[colback=teal!10,colframe=teal!70!black,arc=2mm,boxrule=1pt]
\centering
{\LARGE\bfseries #1}\\[2mm]
{\large Robotique -- \gradelabel{} -- Plan de cours de 2 heures}\\[1mm]
Enseignant : Francis Ki \quad|\quad #2
\end{tcolorbox}}

% Tableau du minutage : 6 contenus (accroche, mini-leçon, guidée, autonome, mise en commun, exit ticket)
\newcommand{\minutage}[6]{%
\begin{tabularx}{\linewidth}{>{\bfseries}p{2.2cm}p{2.6cm}Y p{3cm}}
\toprule
Minutes & Étape & Contenu & Organisation\\
\midrule
0 -- 10 & Accroche & #1 & Classe entière\\
10 -- 30 & Mini-leçon & #2 & Classe entière\\
30 -- 60 & Pratique guidée & #3 & Collectif puis binômes\\
60 -- 95 & Pratique autonome & #4 & Individuel sur poste\\
95 -- 110 & Mise en commun & #5 & Classe entière\\
110 -- 120 & Exit ticket & #6 & Individuel sur poste\\
\bottomrule
\end{tabularx}}

% Étayage progressif : #1 aide 1, #2 aide 2, #3 aide 3
\newcommand{\aides}[3]{%
\textbf{Étayage progressif de l'aide.}
\begin{center}
\begin{tabularx}{\linewidth}{>{\bfseries}p{2.4cm}Y p{4.4cm}}
\toprule
Niveau d'aide & Contenu & Qui intervient\\
\midrule
Aide 1 & #1 & Le binôme seul\\
Aide 2 & #2 & Le binôme avec la carte d'aide\\
Aide 3 & #3 & L'enseignant\\
\bottomrule
\end{tabularx}
\end{center}}

% Les trois groupes de la pratique autonome : #1 tâche A, #2 tâche B, #3 tâche C
\newcommand{\groupes}[3]{%
\textbf{Consigne commune\,:} chaque élève ouvre le fichier de son groupe, le renomme avec son prénom et travaille seul sur son poste. L'enseignant passe d'abord auprès du groupe A pendant 10 min.

\medskip
\begin{tabularx}{\linewidth}{>{\bfseries}p{3cm}Y p{2.2cm}}
\toprule
Groupe & Tâche sur poste & Annexe\\
\midrule
A : en dessous du niveau & #1 & Annexe C\\
B : au niveau & #2 & Annexe D\\
C : au-dessus du niveau & #3 & Annexe E\\
\bottomrule
\end{tabularx}}

% Mise en commun : #1 groupe A, #2 groupe B, #3 groupe C, #4 conclusion
\newcommand{\miseencommun}[4]{%
\begin{etape}{teal}{5. Mise en commun}{95 -- 110 min}
\begin{enumerate}
 \item (2 min) Retour au calme, écrans mis en veille.
 \item (9 min) Trois passages courts de 3 min, avec valorisation :
 \begin{itemize}
  \item \textbf{Groupe A} : #1
  \item \textbf{Groupe B} : #2
  \item \textbf{Groupe C} : #3
 \end{itemize}
 \item (4 min) Conclusion : #4
\end{enumerate}
\end{etape}}

% Clôture : #1 question 1, #2 question 2
\newcommand{\cloture}[2]{%
\begin{etape}{gray}{6. Clôture : exit ticket}{110 -- 120 min}
Formulaire en ligne ou papier (annexe G) :
\begin{enumerate}
 \item #1
 \item #2
 \item Auto-évaluation : \emph{Je maîtrise / Je progresse / J'ai besoin d'aide.}
\end{enumerate}
Lire les réponses avant la séance suivante pour ajuster les groupes de besoin.
\end{etape}}

% Grille d'observation : #1 #2 #3 intitulés de colonnes
\newcommand{\grille}[3]{%
\subsection*{Grille d'observation (enseignant)}
\begin{tabularx}{\linewidth}{Y p{2.4cm} p{2.4cm} p{2.4cm} Y}
\toprule
Élève & #1 & #2 & #3 & Groupe suivant\\
\midrule
 & & & & \\ & & & & \\ & & & & \\ & & & & \\ & & & & \\
\bottomrule
\end{tabularx}}

% Signature de l'exit ticket
\newcommand{\autoeval}{Je me sens : \cb Je maîtrise \quad \cb Je progresse \quad \cb J'ai besoin d'aide}
.
\usepackage[T1]{fontenc}
\usepackage[utf8]{inputenc}
\usepackage{lmodern}
\usepackage[french]{babel}
\usepackage[margin=2cm,headheight=14pt]{geometry}
\usepackage[most]{tcolorbox}
\usepackage{xcolor,tabularx,booktabs,array,enumitem,fancyhdr,amssymb,tikz}
\usetikzlibrary{arrows.meta,positioning}

\pagestyle{fancy}
\fancyhf{}
\lhead{\small Robotique -- \gradelabel{} -- Novembre 2026}
\rhead{\small \leconhdr}
\cfoot{\thepage}
\setlist{itemsep=1pt,topsep=2pt,leftmargin=*}
\renewcommand{\arraystretch}{1.35}
\newcolumntype{Y}{>{\raggedright\arraybackslash}X}
\newcommand{\cb}{$\square$\,}
\newcommand{\fl}{$\rightarrow$}
\newcommand{\lig}{\noindent\rule{\linewidth}{0.3pt}\\[5mm]}

\newtcolorbox{etape}[3]{enhanced,breakable,colback=#1!6,colframe=#1!70!black,
  coltitle=white,fonttitle=\bfseries,title={#2\hfill #3},boxrule=0.8pt,arc=2mm,
  before skip=8pt,after skip=8pt}
\newtcolorbox{astuce}[1][Différenciation et besoins particuliers]{breakable,
  colback=yellow!12,colframe=yellow!50!black,fonttitle=\bfseries\small,
  title={#1},boxrule=0.5pt,arc=1mm,before skip=6pt,after skip=6pt}
\newtcolorbox{code}[1][Code]{breakable,colback=black!4,colframe=black!60,
  fonttitle=\bfseries\small,title={#1},boxrule=0.5pt,arc=1mm,fontupper=\ttfamily\small,
  before skip=6pt,after skip=6pt}
% Blocs de code coloré (écrire le code sans accents : comme dans l'éditeur)
\newtcblisting{arduino}[1][Code Arduino]{listing only,breakable,colback=black!4,
  colframe=black!60,fonttitle=\bfseries\small,title={#1},boxrule=0.5pt,arc=1mm,
  before skip=6pt,after skip=6pt,
  listing options={language=C++,basicstyle=\ttfamily\small,columns=fullflexible,
  keepspaces=true,showstringspaces=false,commentstyle=\color{green!40!black}\itshape,
  keywordstyle=\color{blue!70!black}\bfseries}}
\newtcblisting{python}[1][Code Python]{listing only,breakable,colback=black!4,
  colframe=black!60,fonttitle=\bfseries\small,title={#1},boxrule=0.5pt,arc=1mm,
  before skip=6pt,after skip=6pt,
  listing options={language=Python,basicstyle=\ttfamily\small,columns=fullflexible,
  keepspaces=true,showstringspaces=false,commentstyle=\color{green!40!black}\itshape,
  keywordstyle=\color{blue!70!black}\bfseries}}
\newcommand{\entete}[2]{%
  \begin{tcolorbox}[colback=#1!10,colframe=#1!70!black,boxrule=0.8pt,arc=2mm]
  {\Large\bfseries #2}\hfill Nom : \rule{4cm}{0.3pt}\quad Date : \rule{2.2cm}{0.3pt}
  \end{tcolorbox}}

% Bandeau de titre : #1 titre, #2 semaine
\newcommand{\titreblock}[2]{%
\begin{tcolorbox}[colback=teal!10,colframe=teal!70!black,arc=2mm,boxrule=1pt]
\centering
{\LARGE\bfseries #1}\\[2mm]
{\large Robotique -- \gradelabel{} -- Plan de cours de 2 heures}\\[1mm]
Enseignant : Francis Ki \quad|\quad #2
\end{tcolorbox}}

% Tableau du minutage : 6 contenus (accroche, mini-leçon, guidée, autonome, mise en commun, exit ticket)
\newcommand{\minutage}[6]{%
\begin{tabularx}{\linewidth}{>{\bfseries}p{2.2cm}p{2.6cm}Y p{3cm}}
\toprule
Minutes & Étape & Contenu & Organisation\\
\midrule
0 -- 10 & Accroche & #1 & Classe entière\\
10 -- 30 & Mini-leçon & #2 & Classe entière\\
30 -- 60 & Pratique guidée & #3 & Collectif puis binômes\\
60 -- 95 & Pratique autonome & #4 & Individuel sur poste\\
95 -- 110 & Mise en commun & #5 & Classe entière\\
110 -- 120 & Exit ticket & #6 & Individuel sur poste\\
\bottomrule
\end{tabularx}}

% Étayage progressif : #1 aide 1, #2 aide 2, #3 aide 3
\newcommand{\aides}[3]{%
\textbf{Étayage progressif de l'aide.}
\begin{center}
\begin{tabularx}{\linewidth}{>{\bfseries}p{2.4cm}Y p{4.4cm}}
\toprule
Niveau d'aide & Contenu & Qui intervient\\
\midrule
Aide 1 & #1 & Le binôme seul\\
Aide 2 & #2 & Le binôme avec la carte d'aide\\
Aide 3 & #3 & L'enseignant\\
\bottomrule
\end{tabularx}
\end{center}}

% Les trois groupes de la pratique autonome : #1 tâche A, #2 tâche B, #3 tâche C
\newcommand{\groupes}[3]{%
\textbf{Consigne commune\,:} chaque élève ouvre le fichier de son groupe, le renomme avec son prénom et travaille seul sur son poste. L'enseignant passe d'abord auprès du groupe A pendant 10 min.

\medskip
\begin{tabularx}{\linewidth}{>{\bfseries}p{3cm}Y p{2.2cm}}
\toprule
Groupe & Tâche sur poste & Annexe\\
\midrule
A : en dessous du niveau & #1 & Annexe C\\
B : au niveau & #2 & Annexe D\\
C : au-dessus du niveau & #3 & Annexe E\\
\bottomrule
\end{tabularx}}

% Mise en commun : #1 groupe A, #2 groupe B, #3 groupe C, #4 conclusion
\newcommand{\miseencommun}[4]{%
\begin{etape}{teal}{5. Mise en commun}{95 -- 110 min}
\begin{enumerate}
 \item (2 min) Retour au calme, écrans mis en veille.
 \item (9 min) Trois passages courts de 3 min, avec valorisation :
 \begin{itemize}
  \item \textbf{Groupe A} : #1
  \item \textbf{Groupe B} : #2
  \item \textbf{Groupe C} : #3
 \end{itemize}
 \item (4 min) Conclusion : #4
\end{enumerate}
\end{etape}}

% Clôture : #1 question 1, #2 question 2
\newcommand{\cloture}[2]{%
\begin{etape}{gray}{6. Clôture : exit ticket}{110 -- 120 min}
Formulaire en ligne ou papier (annexe G) :
\begin{enumerate}
 \item #1
 \item #2
 \item Auto-évaluation : \emph{Je maîtrise / Je progresse / J'ai besoin d'aide.}
\end{enumerate}
Lire les réponses avant la séance suivante pour ajuster les groupes de besoin.
\end{etape}}

% Grille d'observation : #1 #2 #3 intitulés de colonnes
\newcommand{\grille}[3]{%
\subsection*{Grille d'observation (enseignant)}
\begin{tabularx}{\linewidth}{Y p{2.4cm} p{2.4cm} p{2.4cm} Y}
\toprule
Élève & #1 & #2 & #3 & Groupe suivant\\
\midrule
 & & & & \\ & & & & \\ & & & & \\ & & & & \\ & & & & \\
\bottomrule
\end{tabularx}}

% Signature de l'exit ticket
\newcommand{\autoeval}{Je me sens : \cb Je maîtrise \quad \cb Je progresse \quad \cb J'ai besoin d'aide}
.
\usepackage[T1]{fontenc}
\usepackage[utf8]{inputenc}
\usepackage{lmodern}
\usepackage[french]{babel}
\usepackage[margin=2cm,headheight=14pt]{geometry}
\usepackage[most]{tcolorbox}
\usepackage{xcolor,tabularx,booktabs,array,enumitem,fancyhdr,amssymb,tikz}
\usetikzlibrary{arrows.meta,positioning}

\pagestyle{fancy}
\fancyhf{}
\lhead{\small Robotique -- \gradelabel{} -- Novembre 2026}
\rhead{\small \leconhdr}
\cfoot{\thepage}
\setlist{itemsep=1pt,topsep=2pt,leftmargin=*}
\renewcommand{\arraystretch}{1.35}
\newcolumntype{Y}{>{\raggedright\arraybackslash}X}
\newcommand{\cb}{$\square$\,}
\newcommand{\fl}{$\rightarrow$}
\newcommand{\lig}{\noindent\rule{\linewidth}{0.3pt}\\[5mm]}

\newtcolorbox{etape}[3]{enhanced,breakable,colback=#1!6,colframe=#1!70!black,
  coltitle=white,fonttitle=\bfseries,title={#2\hfill #3},boxrule=0.8pt,arc=2mm,
  before skip=8pt,after skip=8pt}
\newtcolorbox{astuce}[1][Différenciation et besoins particuliers]{breakable,
  colback=yellow!12,colframe=yellow!50!black,fonttitle=\bfseries\small,
  title={#1},boxrule=0.5pt,arc=1mm,before skip=6pt,after skip=6pt}
\newtcolorbox{code}[1][Code]{breakable,colback=black!4,colframe=black!60,
  fonttitle=\bfseries\small,title={#1},boxrule=0.5pt,arc=1mm,fontupper=\ttfamily\small,
  before skip=6pt,after skip=6pt}
% Blocs de code coloré (écrire le code sans accents : comme dans l'éditeur)
\newtcblisting{arduino}[1][Code Arduino]{listing only,breakable,colback=black!4,
  colframe=black!60,fonttitle=\bfseries\small,title={#1},boxrule=0.5pt,arc=1mm,
  before skip=6pt,after skip=6pt,
  listing options={language=C++,basicstyle=\ttfamily\small,columns=fullflexible,
  keepspaces=true,showstringspaces=false,commentstyle=\color{green!40!black}\itshape,
  keywordstyle=\color{blue!70!black}\bfseries}}
\newtcblisting{python}[1][Code Python]{listing only,breakable,colback=black!4,
  colframe=black!60,fonttitle=\bfseries\small,title={#1},boxrule=0.5pt,arc=1mm,
  before skip=6pt,after skip=6pt,
  listing options={language=Python,basicstyle=\ttfamily\small,columns=fullflexible,
  keepspaces=true,showstringspaces=false,commentstyle=\color{green!40!black}\itshape,
  keywordstyle=\color{blue!70!black}\bfseries}}
\newcommand{\entete}[2]{%
  \begin{tcolorbox}[colback=#1!10,colframe=#1!70!black,boxrule=0.8pt,arc=2mm]
  {\Large\bfseries #2}\hfill Nom : \rule{4cm}{0.3pt}\quad Date : \rule{2.2cm}{0.3pt}
  \end{tcolorbox}}

% Bandeau de titre : #1 titre, #2 semaine
\newcommand{\titreblock}[2]{%
\begin{tcolorbox}[colback=teal!10,colframe=teal!70!black,arc=2mm,boxrule=1pt]
\centering
{\LARGE\bfseries #1}\\[2mm]
{\large Robotique -- \gradelabel{} -- Plan de cours de 2 heures}\\[1mm]
Enseignant : Francis Ki \quad|\quad #2
\end{tcolorbox}}

% Tableau du minutage : 6 contenus (accroche, mini-leçon, guidée, autonome, mise en commun, exit ticket)
\newcommand{\minutage}[6]{%
\begin{tabularx}{\linewidth}{>{\bfseries}p{2.2cm}p{2.6cm}Y p{3cm}}
\toprule
Minutes & Étape & Contenu & Organisation\\
\midrule
0 -- 10 & Accroche & #1 & Classe entière\\
10 -- 30 & Mini-leçon & #2 & Classe entière\\
30 -- 60 & Pratique guidée & #3 & Collectif puis binômes\\
60 -- 95 & Pratique autonome & #4 & Individuel sur poste\\
95 -- 110 & Mise en commun & #5 & Classe entière\\
110 -- 120 & Exit ticket & #6 & Individuel sur poste\\
\bottomrule
\end{tabularx}}

% Étayage progressif : #1 aide 1, #2 aide 2, #3 aide 3
\newcommand{\aides}[3]{%
\textbf{Étayage progressif de l'aide.}
\begin{center}
\begin{tabularx}{\linewidth}{>{\bfseries}p{2.4cm}Y p{4.4cm}}
\toprule
Niveau d'aide & Contenu & Qui intervient\\
\midrule
Aide 1 & #1 & Le binôme seul\\
Aide 2 & #2 & Le binôme avec la carte d'aide\\
Aide 3 & #3 & L'enseignant\\
\bottomrule
\end{tabularx}
\end{center}}

% Les trois groupes de la pratique autonome : #1 tâche A, #2 tâche B, #3 tâche C
\newcommand{\groupes}[3]{%
\textbf{Consigne commune\,:} chaque élève ouvre le fichier de son groupe, le renomme avec son prénom et travaille seul sur son poste. L'enseignant passe d'abord auprès du groupe A pendant 10 min.

\medskip
\begin{tabularx}{\linewidth}{>{\bfseries}p{3cm}Y p{2.2cm}}
\toprule
Groupe & Tâche sur poste & Annexe\\
\midrule
A : en dessous du niveau & #1 & Annexe C\\
B : au niveau & #2 & Annexe D\\
C : au-dessus du niveau & #3 & Annexe E\\
\bottomrule
\end{tabularx}}

% Mise en commun : #1 groupe A, #2 groupe B, #3 groupe C, #4 conclusion
\newcommand{\miseencommun}[4]{%
\begin{etape}{teal}{5. Mise en commun}{95 -- 110 min}
\begin{enumerate}
 \item (2 min) Retour au calme, écrans mis en veille.
 \item (9 min) Trois passages courts de 3 min, avec valorisation :
 \begin{itemize}
  \item \textbf{Groupe A} : #1
  \item \textbf{Groupe B} : #2
  \item \textbf{Groupe C} : #3
 \end{itemize}
 \item (4 min) Conclusion : #4
\end{enumerate}
\end{etape}}

% Clôture : #1 question 1, #2 question 2
\newcommand{\cloture}[2]{%
\begin{etape}{gray}{6. Clôture : exit ticket}{110 -- 120 min}
Formulaire en ligne ou papier (annexe G) :
\begin{enumerate}
 \item #1
 \item #2
 \item Auto-évaluation : \emph{Je maîtrise / Je progresse / J'ai besoin d'aide.}
\end{enumerate}
Lire les réponses avant la séance suivante pour ajuster les groupes de besoin.
\end{etape}}

% Grille d'observation : #1 #2 #3 intitulés de colonnes
\newcommand{\grille}[3]{%
\subsection*{Grille d'observation (enseignant)}
\begin{tabularx}{\linewidth}{Y p{2.4cm} p{2.4cm} p{2.4cm} Y}
\toprule
Élève & #1 & #2 & #3 & Groupe suivant\\
\midrule
 & & & & \\ & & & & \\ & & & & \\ & & & & \\ & & & & \\
\bottomrule
\end{tabularx}}

% Signature de l'exit ticket
\newcommand{\autoeval}{Je me sens : \cb Je maîtrise \quad \cb Je progresse \quad \cb J'ai besoin d'aide}


\begin{document}
\titreblock{Simulation d'un feu intelligent}{Semaine du 23 au 27 novembre 2026}

%==================================================================
\section{Informations générales}
\begin{tabularx}{\linewidth}{>{\bfseries}p{3.6cm}Y}
\toprule
Titre & Simulation d'un feu intelligent\\
Classe / effectif & Grade 9, 11 à 20 élèves, niveaux hétérogènes, élèves à besoins particuliers\\
Durée & 2 heures (120 min)\\
Lieu & Salle informatique, un poste par élève\\
Matériel et logiciels &
\begin{itemize}
 \item Un poste par élève avec navigateur web : \textbf{Tinkercad Circuits} (simulation) et son moniteur série.
 \item Arduino IDE sur les postes pour vérifier le code.
 \item Facultatif : cartes Arduino avec LEDs et résistances pour tester un montage réel.
 \item Vidéoprojecteur ou tableau interactif.
\end{itemize}\\
Pré-requis &
Séance précédente : cycle d'états du feu tricolore, tableau des états, temporisations. Variables de type \texttt{int}.\\
\bottomrule
\end{tabularx}

\begin{astuce}[Préparation avant le cours (environ 45 min)]
\begin{enumerate}
 \item Préparer le circuit du feu tricolore de la séance précédente (broches 8, 9, 10) comme base commune.
 \item Préparer une courte vidéo ou des photos de carrefours : carrefour très chargé et carrefour vide la nuit.
 \item Préparer les fichiers des annexes C à E et le formulaire de sortie.
 \item Préciser que les capteurs (bouton, ultrason) seront ajoutés en janvier : ici, le «~trafic~» et l'«~heure~» sont simulés par des variables.
\end{enumerate}
\end{astuce}

%==================================================================
\section{Objectifs d'apprentissage}
\begin{tcolorbox}[colback=gray!8,colframe=gray!60!black,arc=1mm,boxrule=0.5pt]
\textbf{Objectifs du programme}
\begin{itemize}
 \item Analyser un besoin réel.
 \item Réaliser un programme complet.
\end{itemize}
\textbf{À la fin de la séance, l'élève est capable\,:}
\begin{itemize}
 \item de rédiger un mini cahier des charges (besoin, fonctions, scénarios)\,;
 \item d'utiliser des variables et une condition \texttt{if} pour adapter le comportement du feu\,;
 \item de tester chaque scénario en simulation et de documenter ses résultats.
\end{itemize}
\textbf{Compétences numériques\,:} structurer un programme, utiliser le moniteur série pour contrôler, écrire une courte documentation technique.
\end{tcolorbox}

%==================================================================
\section{Déroulé minuté (120 min)}
\minutage{Pourquoi un feu devrait-il être intelligent ?}{Besoin, scénarios et variables}{Un mode nuit en collectif, puis un feu qui prolonge le vert en binômes}{Trois niveaux de besoin}{Valoriser chaque groupe}{Trois questions sur poste}

%------------------------------------------------------------------
\begin{etape}{orange}{1. Accroche : le feu qui s'adapte}{0 -- 10 min}
\begin{enumerate}
 \item (5 min) Projeter deux situations : un carrefour bloqué aux heures de pointe, un feu rouge à 3 h du matin sans voiture. \emph{Qu'est-ce qui ne va pas ? Que faudrait-il changer ?}
 \item (5 min) Noter les idées au tableau : feu plus long quand il y a du trafic, clignotant la nuit. Question-problème : \textbf{«~Comment faire adapter un feu à la situation ?~»}
\end{enumerate}
\end{etape}

%------------------------------------------------------------------
\begin{etape}{blue}{2. Mini-leçon : du besoin au programme}{10 -- 30 min}
\textbf{Support\,:} annexe A. Quatre temps de 5 min.

\textbf{Temps 1 : le besoin.} \emph{Qui est concerné ? Quel problème ? Quelle amélioration ?}

\textbf{Temps 2 : des scénarios.}
\begin{center}
\begin{tabularx}{\linewidth}{>{\bfseries}p{2.8cm}Y Y}
\toprule
Scénario & Condition (variable) & Comportement du feu\\
\midrule
Normal & \texttt{nuit == false} et \texttt{trafic == 0} & Cycle vert 5 s, orange 2 s, rouge 5 s\\
Trafic dense & \texttt{trafic == 1} & Vert prolongé à 8 s\\
Nuit & \texttt{nuit == true} & Orange clignotant (500 ms)\\
\bottomrule
\end{tabularx}
\end{center}

\textbf{Temps 3 : variables et condition.}
\begin{arduino}[Structure avec variables]
bool nuit = false;      // mode nuit ou non
int trafic = 0;         // 0 = faible, 1 = dense
int cycle = 0;          // compte les cycles

void loop() {
  if (nuit) {
    digitalWrite(9, HIGH); delay(500);
    digitalWrite(9, LOW);  delay(500);
  } else {
    // cycle normal, avec vert plus long si trafic dense
  }
}
\end{arduino}

\textbf{Temps 4 : simuler le monde réel.} Sans capteur, on change les variables pour simuler : le programme compte les cycles et passe en mode nuit après 4 cycles, puis revient en mode jour après 4 autres.

\begin{astuce}
Écrire les scénarios sous forme de phrases «~si\ldots{} alors\ldots~» avant de coder. Utiliser des couleurs pour distinguer les variables du reste du code. Lire le code à voix haute avec les noms des variables.
\end{astuce}
\end{etape}

%------------------------------------------------------------------
\begin{etape}{green}{3. Pratique guidée}{30 -- 60 min}
\textbf{Fiche\,:} annexe B.

\textbf{3a. Le mode nuit en collectif (12 min).} Ajouter à la base commune une variable \texttt{nuit} et la condition \texttt{if (nuit)}. Tester avec \texttt{nuit = true} puis \texttt{false}.

\textbf{3b. Le vert prolongé en binômes (18 min).} Ajouter la variable \texttt{trafic} : si \texttt{trafic == 1}, le vert dure 8 s au lieu de 5 s. Tester les deux valeurs. Alternance des rôles.

\aides{Le programme du 3a et le tableau des scénarios}{Questions guides : \emph{Quelle variable change le comportement ? Quelle valeur de \texttt{delay} dans chaque cas ?}}{L'enseignant écrit la structure \texttt{if / else} avec le binôme, qui complète}

\begin{astuce}
Tester un scénario à la fois en changeant seulement la valeur d'une variable. Commenter chaque \texttt{if} avec le scénario correspondant.
\end{astuce}
\end{etape}

%------------------------------------------------------------------
\begin{etape}{violet}{4. Pratique autonome différenciée}{60 -- 95 min}
\groupes{Programme du mode nuit à compléter (conditions et valeurs à placer avec une banque de mots), puis test avec \texttt{nuit = true} et \texttt{false} et relevé des résultats dans un tableau.}{Réaliser un feu avec deux scénarios (nuit et trafic dense) et un compteur de cycles qui bascule en mode nuit après 4 cycles. Tester chaque scénario et noter les résultats.}{Concevoir un feu intelligent complet : cahier des charges d'une page, trois scénarios, code structuré avec des fonctions, affichage du scénario dans le moniteur série, et un tableau de tests.}

\medskip
\textbf{Plan B si un élève termine tôt\,:} ajouter un feu piéton qui s'adapte aux mêmes scénarios.
\end{etape}

%------------------------------------------------------------------
\miseencommun{un élève montre son mode nuit et explique la condition \texttt{if}. Valoriser le lien entre le scénario et le code.}{un binôme montre le basculement jour et nuit par compteur de cycles. Valoriser le test méthodique.}{un élève présente son cahier des charges et son tableau de tests. Valoriser la rigueur de la démarche.}{\textit{un feu «~intelligent~» adapte son comportement avec des variables et des conditions. Dans le trimestre suivant, on le fera avec de vrais capteurs (boutons, ultrason).} Annoncer la préparation du projet pour décembre.}

\cloture{À quoi sert la variable \texttt{nuit} dans le programme du feu ?}{Écris la condition qui fait durer le vert 8 s quand le trafic est dense.}

%==================================================================
% ANNEXES
%==================================================================
\newpage
\section*{Annexe A : Support résumé de la mini-leçon}
\entete{blue}{Du besoin au programme}

\begin{tcolorbox}[colback=yellow!10,colframe=yellow!50!black]
\textbf{Besoin} $\rightarrow$ \textbf{scénarios} $\rightarrow$ \textbf{variables et conditions} $\rightarrow$ \textbf{programme} $\rightarrow$ \textbf{tests}
\end{tcolorbox}

\begin{tabularx}{\linewidth}{>{\bfseries}p{2.8cm}Y Y}
\toprule
Scénario & Condition & Comportement\\
\midrule
Normal & ni nuit ni trafic dense & Cycle 5 s / 2 s / 5 s\\
Trafic dense & \texttt{trafic == 1} & Vert 8 s\\
Nuit & \texttt{nuit == true} & Orange clignotant\\
\bottomrule
\end{tabularx}

\begin{arduino}[Modèle]
if (nuit) {
  // orange clignotant
} else {
  if (trafic == 1) { delay(8000); }
  else             { delay(5000); }
}
\end{arduino}

%==================================================================
\newpage
\section*{Annexe B : Fiche de pratique guidée}
\entete{green}{Pratique guidée : le feu intelligent}

\subsection*{3a. Mode nuit}
Variable utilisée : \dotfill\quad Valeur pour la nuit : \dotfill\\
Comportement observé : \dotfill

\subsection*{3b. Tableau de tests}
\begin{tabularx}{\linewidth}{>{\bfseries}p{2.4cm}Y Y Y}
\toprule
Test & Valeur des variables & Comportement attendu & Observé ?\\
\midrule
1 & \texttt{nuit = false}, \texttt{trafic = 0} & & \\
2 & \texttt{nuit = false}, \texttt{trafic = 1} & & \\
3 & \texttt{nuit = true} & & \\
\bottomrule
\end{tabularx}

\begin{tcolorbox}[colback=yellow!10,colframe=yellow!50!black,title={Cartes d'aide à utiliser dans l'ordre},fonttitle=\bfseries\small]
\textbf{Aide 1 :} regarde le programme du mode nuit et le tableau des scénarios.\\
\textbf{Aide 2 :} Quelle variable change le comportement ? Quelle valeur de \texttt{delay} dans chaque cas ?\\
\textbf{Aide 3 :} demande à l'enseignant d'écrire le \texttt{if / else} avec toi.
\end{tcolorbox}

%==================================================================
\newpage
\section*{Annexe C : Niveau A (en dessous du niveau)}
\entete{orange}{Niveau A : je complète le mode nuit}

\textbf{Banque :} nuit, true, false, HIGH, LOW, 500.

\begin{arduino}[Programme à compléter]
bool nuit = _____;

void loop() {
  if (_____) {
    digitalWrite(9, HIGH); delay(_____);
    digitalWrite(9, _____); delay(500);
  } else {
    // cycle normal
  }
}
\end{arduino}

\subsection*{Activité : relevé des tests}
\begin{tabularx}{\linewidth}{>{\bfseries}p{3.4cm}Y Y}
\toprule
Valeur de \texttt{nuit} & Ce que fait le feu & Conforme ?\\
\midrule
\texttt{true} & & \\ \texttt{false} & & \\
\bottomrule
\end{tabularx}

%==================================================================
\newpage
\section*{Annexe D : Niveau B (au niveau)}
\entete{blue}{Niveau B : deux scénarios et un compteur}

\textbf{Mission :} le feu passe en mode nuit après 4 cycles, puis en mode jour après 4 autres. En mode jour, le vert dure 8 s si \texttt{trafic == 1}.

\begin{tabularx}{\linewidth}{>{\bfseries}p{2.8cm}Y Y Y}
\toprule
Cycle n° & Mode & Durée du vert & Observé ?\\
\midrule
1 & & & \\ 2 & & & \\ 3 & & & \\ 4 & & & \\ 5 & & & \\
\bottomrule
\end{tabularx}

\subsection*{Questions}
\begin{enumerate}
 \item Quelle ligne remet \texttt{cycle} à zéro ? \dotfill
 \item Pourquoi tester chaque scénario séparément ? \dotfill
\end{enumerate}

%==================================================================
\newpage
\section*{Annexe E : Niveau C (au-dessus du niveau)}
\entete{violet}{Niveau C : mon feu intelligent complet}

\subsection*{Étape 1 : mini cahier des charges}
Besoin : \dotfill\\
Utilisateurs : \dotfill\\
Trois scénarios : \lig

\subsection*{Étape 2 : tableau de tests}
\begin{tabularx}{\linewidth}{>{\bfseries}p{2.4cm}Y Y Y}
\toprule
Scénario & Valeur des variables & Résultat attendu & Résultat observé\\
\midrule
 & & & \\ & & & \\ & & & \\
\bottomrule
\end{tabularx}

\subsection*{Étape 3 : moniteur série}
Affiche le scénario avec \texttt{Serial.println("nuit");}. Recopie ici une ligne de ton moniteur : \dotfill

\textbf{Défi bonus :} ajoute une variable \texttt{priorite} pour donner la priorité à une voie.

%==================================================================
\newpage
\section*{Annexe F : Corrections et attendus (enseignant)}

\subsection*{Annexe B}
Tests : 1 cycle normal (vert 5 s) ; 2 vert prolongé (8 s) ; 3 orange clignotant.

\subsection*{Annexe C}
\texttt{false}, \texttt{nuit}, \texttt{500}, \texttt{LOW}. \texttt{true} $\rightarrow$ orange clignotant ; \texttt{false} $\rightarrow$ cycle normal.

\subsection*{Annexe D}
Cycles 1 à 4 : jour ; 5 à 8 : nuit. Remise à zéro de \texttt{cycle} : \texttt{cycle = 0;} à la bascule. On teste séparément pour localiser l'erreur.

\subsection*{Annexe E}
Un bon dossier contient un besoin clair, trois scénarios cohérents, un tableau de tests avec résultat observé, et un code lisible avec fonctions.\\
\textbf{Critères de réussite :} chaque scénario testé et conforme, code commenté, moniteur série utilisé.

%==================================================================
\newpage
\section*{Annexe G : Exit ticket et grille d'observation}
\entete{gray}{Exit ticket : feu intelligent}

\begin{enumerate}
 \item À quoi sert la variable \texttt{nuit} dans le programme du feu ?\\ \lig
 \item Écris la condition qui fait durer le vert 8 s quand le trafic est dense.\\ \lig
 \item \autoeval
\end{enumerate}
\textbf{Attendus :} 1. À choisir entre le cycle normal et l'orange clignotant. 2. \texttt{if (trafic == 1) \{ delay(8000); \}}.

\bigskip
\grille{Scénarios}{Tests réalisés}{Code lisible}

\end{document}
